\chapter{Complexity — Straight Line}
%%% generated by the blueprint pipeline from TCSlib.Complexity.CircuitComplexity.StraightLine
\section{Overview}

Declarations of the Lean module \texttt{TCSlib.Complexity.CircuitComplexity.StraightLine}.

\section{Declarations}

\begin{definition}[Straight-line program operands]
\lean{BoolCircuit.SLOperand}
\label{BoolCircuit.SLOperand}
\leanok
For each natural number $n$, an operand is either an input variable $x_k$ for some
$k	ext{ in Fin }n$, its negation, or the value of a numbered program line.
\end{definition}

\begin{definition}[Straight-line operations]
\lean{BoolCircuit.SLOp}
\label{BoolCircuit.SLOp}
\leanok
A straight-line operation is either conjunction or disjunction.
\end{definition}

\begin{definition}[Straight-line statement]
\lean{BoolCircuit.SLStmt}
\label{BoolCircuit.SLStmt}
\leanok
\uses{BoolCircuit.SLOp, BoolCircuit.SLOperand}
For each natural number $n$, a straight-line statement consists of an operation, either
conjunction or disjunction, together with two operands drawn from the input variables,
their negations, or program-line values.
\end{definition}

\begin{definition}[Boolean operation evaluation]
\lean{BoolCircuit.SLOp.apply}
\label{BoolCircuit.SLOp.apply}
\leanok
\uses{BoolCircuit.SLOp}
Define the evaluation of a straight-line operation on two Boolean values: conjunction
returns their conjunction, and disjunction returns their disjunction.
\end{definition}

\begin{definition}[Dual Boolean operation]
\lean{BoolCircuit.SLOp.dual}
\label{BoolCircuit.SLOp.dual}
\leanok
\uses{BoolCircuit.SLOp}
The dual of conjunction is disjunction, and the dual of disjunction is conjunction.
\end{definition}

\begin{definition}[Gate kind of an operation]
\lean{BoolCircuit.SLOp.kind}
\label{BoolCircuit.SLOp.kind}
\leanok
\uses{BoolCircuit.GateKind, BoolCircuit.SLOp}
The gate kind of a straight-line operation is the corresponding Boolean gate:
conjunction has kind AND, and disjunction has kind OR.
\end{definition}

\begin{definition}[Operand precedes a line]
\lean{BoolCircuit.SLOperand.Before}
\label{BoolCircuit.SLOperand.Before}
\leanok
\uses{BoolCircuit.SLOperand}
An operand precedes line $i$ if it is a previously computed line with index strictly
less than $i$; every input and negated input precedes every line.
\end{definition}

\begin{definition}[Evaluation of an operand]
\lean{BoolCircuit.SLOperand.eval}
\label{BoolCircuit.SLOperand.eval}
\leanok
\uses{BoolCircuit.SLOperand}
Given an assignment of Boolean values to the $n$ inputs and a list of previously
computed Boolean values, an input operand evaluates to its assigned value, a negated
input operand to its complement, and a line operand to the list entry at its index, or
to false if that index is out of range.
\end{definition}

\begin{theorem}[Monotonicity of the before relation]
\lean{BoolCircuit.SLOperand.Before.mono}
\label{BoolCircuit.SLOperand.Before.mono}
\leanok
\uses{BoolCircuit.SLOperand, BoolCircuit.SLOperand.Before}
\difficulty{4}
Let $o$ be a straight-line operand and let $i,j$ be natural numbers. If $o$ is before
line $i$ and $i\le j$, then $o$ is before line $j$.
\end{theorem}

\begin{theorem}[Operand evaluation under list extension]
\lean{BoolCircuit.SLOperand.eval_append}
\label{BoolCircuit.SLOperand.eval_append}
\leanok
\uses{BoolCircuit.SLOperand, BoolCircuit.SLOperand.Before, BoolCircuit.SLOperand.eval}
\difficulty{4}
Let $x:	ext{Fin}(n)	o	ext{Bool}$ be an assignment, let $o$ be an operand, and let $ys$
and $ext$ be Boolean lists. If every line operand referenced by $o$ has index less than
the length of $ys$ (with input and negated-input operands satisfying this condition
automatically), then evaluating $o$ on $ys$ appended with $ext$ gives the same value as
evaluating it on $ys$.
\end{theorem}

\begin{definition}[Evaluation of a statement]
\lean{BoolCircuit.SLStmt.eval}
\label{BoolCircuit.SLStmt.eval}
\leanok
\uses{BoolCircuit.SLOp.apply, BoolCircuit.SLOperand.eval, BoolCircuit.SLStmt}
A statement evaluates by applying its Boolean operation to the evaluations of its two
operands, using the given input assignment and list of previously computed values.
\end{definition}

\begin{definition}[Statement precedes a line]
\lean{BoolCircuit.SLStmt.Before}
\label{BoolCircuit.SLStmt.Before}
\leanok
\uses{BoolCircuit.SLOperand.Before, BoolCircuit.SLStmt}
A statement precedes line $i$ exactly when each of its two operands precedes line $i$.
\end{definition}

\begin{definition}[Validity of a straight-line program]
\lean{BoolCircuit.StraightLine.LinesValid}
\label{BoolCircuit.StraightLine.LinesValid}
\leanok
\uses{BoolCircuit.SLStmt, BoolCircuit.SLStmt.Before}
A list of Boolean straight-line statements is valid if, at every position $i$ in the
list, each operand is either an input variable or its negation, or the value computed at
some earlier position $j<i$.
\end{definition}

\begin{theorem}[Validity of the empty list]
\lean{BoolCircuit.StraightLine.LinesValid.nil}
\label{BoolCircuit.StraightLine.LinesValid.nil}
\leanok
\uses{BoolCircuit.SLStmt, BoolCircuit.StraightLine.LinesValid}
\difficulty{1}
For every number $n$ of input variables, the empty list of straight-line statements is
valid: each statement in the list uses only input variables or earlier lines.
\end{theorem}

\begin{theorem}[Validity after appending a statement]
\lean{BoolCircuit.StraightLine.LinesValid.snoc}
\label{BoolCircuit.StraightLine.LinesValid.snoc}
\leanok
\uses{BoolCircuit.SLStmt, BoolCircuit.SLStmt.Before, BoolCircuit.StraightLine.LinesValid}
\difficulty{4}
Let $n$ be the number of input variables, and let $ss$ be a valid list of straight-line
statements. If a statement uses only input variables and lines with index less than the
length of $ss$, then appending it to $ss$ gives another valid list.
\end{theorem}

\begin{definition}[Execution of straight-line statements]
\lean{BoolCircuit.StraightLine.runLines}
\label{BoolCircuit.StraightLine.runLines}
\leanok
\uses{BoolCircuit.SLStmt, BoolCircuit.SLStmt.eval}
Let $n$ be a natural number, let $ss$ be a finite list of straight-line statements on
$n$ inputs, and let $x$ assign a Boolean value to each input. The list of line values
produced by running $ss$ on $x$ is obtained by starting with no line values and, in
order, appending the Boolean value of each statement evaluated at $x$ using the line
values already produced.
\end{definition}

\begin{theorem}[Empty Straight-Line Program]
\lean{BoolCircuit.StraightLine.runLines_nil}
\label{BoolCircuit.StraightLine.runLines_nil}
\leanok
\uses{BoolCircuit.SLStmt, BoolCircuit.StraightLine.runLines}
\difficulty{1}
For any natural number $n$ and any Boolean assignment $x : 	ext{Fin}(n) \to
\{\mathrm{false},\mathrm{true}\}$, running the empty list of straight-line statements on
$x$ returns the empty list.
\end{theorem}

\begin{theorem}[Evaluation After Appending One Statement]
\lean{BoolCircuit.StraightLine.runLines_snoc}
\label{BoolCircuit.StraightLine.runLines_snoc}
\leanok
\uses{BoolCircuit.SLStmt, BoolCircuit.SLStmt.eval, BoolCircuit.StraightLine.runLines}
\difficulty{3}
For any natural number $n$, any list $ss$ of straight-line statements on $n$ Boolean
inputs, any statement $s$ on those inputs, and any Boolean input assignment $x$,
evaluating $ss$ followed by $s$ gives the evaluation of $ss$ with the value of $s$ on
that evaluation appended.
\end{theorem}

\begin{theorem}[Length of the Evaluated Lines]
\lean{BoolCircuit.StraightLine.length_runLines}
\label{BoolCircuit.StraightLine.length_runLines}
\leanok
\uses{BoolCircuit.SLStmt, BoolCircuit.StraightLine.runLines, BoolCircuit.StraightLine.runLines_snoc}
\difficulty{4}
For any natural number $n$, any finite list $ss$ of straight-line statements on $n$
Boolean inputs, and any Boolean input assignment $x$, evaluating the statements in order
produces a list whose length is $|ss|$.
\end{theorem}

\begin{definition}[Straight-line Boolean program]
\lean{BoolCircuit.StraightLineProgram}
\label{BoolCircuit.StraightLineProgram}
\leanok
\uses{BoolCircuit.SLStmt, BoolCircuit.StraightLine.LinesValid}
For a natural number $n$, a straight-line Boolean program on $n$ input bits is a finite
ordered list of statements such that each statement reads only input variables or values
computed by earlier statements.
\end{definition}

\begin{definition}[Program length]
\lean{BoolCircuit.StraightLineProgram.length}
\label{BoolCircuit.StraightLineProgram.length}
\leanok
For a straight-line Boolean program $P$, its length is the number of statements in its
ordered list.
\end{definition}

\begin{definition}[Program line values]
\lean{BoolCircuit.StraightLineProgram.values}
\label{BoolCircuit.StraightLineProgram.values}
\leanok
\uses{BoolCircuit.StraightLine.runLines}
Given a straight-line Boolean program $P$ on $n$ input bits and an input
$x\\in\\{0,1\\}^n$, the list of values produced by its statements is the result of
executing those statements in order on $x$.
\end{definition}

\begin{definition}[Program evaluation]
\lean{BoolCircuit.StraightLineProgram.eval}
\label{BoolCircuit.StraightLineProgram.eval}
\leanok
\uses{BoolCircuit.StraightLineProgram.length, BoolCircuit.StraightLineProgram.values}
Given a straight-line Boolean program $P$ on $n$ input bits and an input
$x\\in\\{0,1\\}^n$, its output is the value at position $|P|-1$ in the list of statement
values, using $0$ if that position is unavailable.
\end{definition}

\begin{theorem}[Evaluation with No Inputs]
\lean{BoolCircuit.StraightLineProgram.eval_of_zero}
\label{BoolCircuit.StraightLineProgram.eval_of_zero}
\leanok
\uses{BoolCircuit.StraightLineProgram, BoolCircuit.StraightLineProgram.eval, BoolCircuit.StraightLineProgram.length, BoolCircuit.StraightLineProgram.values}
\difficulty{5}
Let $P$ be a straight-line program with zero input variables, and let $x$ be an
assignment of Boolean values to its input variables. Then the evaluation of $P$ at $x$
is false.
\end{theorem}

\begin{definition}[Operand vertex index]
\lean{BoolCircuit.SLOperand.vertex}
\label{BoolCircuit.SLOperand.vertex}
\leanok
\uses{BoolCircuit.SLOperand}
For a straight-line operand with $n$ input variables, assign vertex index $k$ to input
$x_k$, index $n+k$ to its negation, and index $2n+j$ to the value computed on line $j$.
\end{definition}

\begin{definition}[Statement as a gate]
\lean{BoolCircuit.SLStmt.toGate}
\label{BoolCircuit.SLStmt.toGate}
\leanok
\uses{BoolCircuit.DAGGate, BoolCircuit.SLOp.kind, BoolCircuit.SLOperand.vertex, BoolCircuit.SLStmt}
Represent a straight-line statement by a gate labeled with its conjunction or
disjunction operation. Its argument list contains the vertex indices of its two
operands, with the repeated index listed only once when both operands have the same
index.
\end{definition}

\begin{definition}[Negation gates]
\lean{BoolCircuit.StraightLine.negGates}
\label{BoolCircuit.StraightLine.negGates}
\leanok
\uses{BoolCircuit.DAGGate}
For $n$ input variables, form the list of $n$ negation gates whose respective single
arguments are the indices $0,1,\ldots,n-1$.
\end{definition}

\begin{definition}[Straight-line program gates]
\lean{BoolCircuit.StraightLine.slpGates}
\label{BoolCircuit.StraightLine.slpGates}
\leanok
\uses{BoolCircuit.DAGGate, BoolCircuit.SLStmt, BoolCircuit.SLStmt.toGate, BoolCircuit.StraightLine.negGates}
For a list of straight-line statements on $n$ input variables, form the gate list by
concatenating the $n$ input negation gates with the gate representations of the
statements in their given order.
\end{definition}

\begin{theorem}[Length of the negation-gate list]
\lean{BoolCircuit.StraightLine.length_negGates}
\label{BoolCircuit.StraightLine.length_negGates}
\leanok
\uses{BoolCircuit.StraightLine.negGates}
\difficulty{1}
For every natural number $n$, the list consisting of one unary negation gate for each
index from $0$ through $n-1$ has length $n$.
\end{theorem}

\begin{theorem}[Length of the Gate List]
\lean{BoolCircuit.StraightLine.length_slpGates}
\label{BoolCircuit.StraightLine.length_slpGates}
\leanok
\uses{BoolCircuit.SLStmt, BoolCircuit.StraightLine.slpGates}
\difficulty{1}
For any natural number $n$ and any finite list $ss$ of straight-line statements on $n$
inputs, the gate list formed from the $n$ input-negation gates followed by one gate for
each statement has length $n+|ss|$.
\end{theorem}

\begin{theorem}[Appending a straight-line statement]
\lean{BoolCircuit.StraightLine.slpGates_snoc}
\label{BoolCircuit.StraightLine.slpGates_snoc}
\leanok
\uses{BoolCircuit.SLStmt, BoolCircuit.SLStmt.toGate, BoolCircuit.StraightLine.slpGates}
\difficulty{1}
For any natural number $n$, any finite list $ss$ of straight-line statements on $n$
inputs, and any such statement $s$, the gate list for $ss$ followed by $s$ is the gate
list for $ss$ with the gate representing $s$ appended.
\end{theorem}

\begin{theorem}[Evaluation of a statement gate]
\lean{BoolCircuit.SLStmt.toGate_eval}
\label{BoolCircuit.SLStmt.toGate_eval}
\leanok
\uses{BoolCircuit.DAGGate.eval, BoolCircuit.SLOp.apply, BoolCircuit.SLOp.kind, BoolCircuit.SLOperand.vertex, BoolCircuit.SLStmt, BoolCircuit.SLStmt.toGate}
\difficulty{4}
Let $s$ be a straight-line statement, with operation $ullet$ (conjunction or
disjunction) and operands $a$ and $b$, and let $vals$ be a list of Boolean vertex
values. The gate associated with $s$ evaluates to $ullet(vals[v_a],vals[v_b])$, where
$v_a$ and $v_b$ are the vertices associated with its two operands and any out-of-range
list lookup is taken to be false.
\end{theorem}

\begin{theorem}[Acyclicity of the Negation Gates]
\lean{BoolCircuit.StraightLine.negGates_acyclic}
\label{BoolCircuit.StraightLine.negGates_acyclic}
\leanok
\uses{BoolCircuit.GatesAcyclic, BoolCircuit.StraightLine.length_negGates, BoolCircuit.StraightLine.negGates}
\difficulty{3}
For every natural number $n$, the list of $n$ gates that negate the input vertices
$0,1,\\ldots,n-1$ is acyclic: for each gate at position $i$ and each vertex it reads,
that vertex has index less than $n+i$.
\end{theorem}

\begin{theorem}[Vertex bound for an operand]
\lean{BoolCircuit.SLOperand.vertex_lt}
\label{BoolCircuit.SLOperand.vertex_lt}
\leanok
\uses{BoolCircuit.SLOperand, BoolCircuit.SLOperand.Before, BoolCircuit.SLOperand.vertex}
\difficulty{4}
Let $o$ be an operand and $i$ a natural number. If every line operand referenced by $o$
has index less than $i$, then the vertex associated with $o$ is less than $2n+i$.
\end{theorem}

\begin{theorem}[Acyclicity of the straight-line gate list]
\lean{BoolCircuit.StraightLine.slpGates_acyclic}
\label{BoolCircuit.StraightLine.slpGates_acyclic}
\leanok
\uses{BoolCircuit.GatesAcyclic, BoolCircuit.GatesAcyclic.snoc, BoolCircuit.SLOperand.vertex_lt, BoolCircuit.SLStmt, BoolCircuit.SLStmt.Before, BoolCircuit.SLStmt.toGate, BoolCircuit.StraightLine.LinesValid, BoolCircuit.StraightLine.negGates_acyclic, BoolCircuit.StraightLine.slpGates, BoolCircuit.StraightLine.slpGates_snoc}
\difficulty{5}
Let $n$ be a natural number, and let $ss$ be a finite list of straight-line statements
on $n$ inputs. Assume that at every position $i<|ss|$, each line operand in the
statement at position $i$ refers to a line with index less than $i$. Then, in the list
formed by the $n$ input-negation gates followed by the gates corresponding to the
statements in $ss$, every gate at position $i$ reads only vertices with indices less
than $n+i$.
\end{theorem}

\begin{theorem}[Negated input gate values]
\lean{BoolCircuit.StraightLine.vertexValue_negGates}
\label{BoolCircuit.StraightLine.vertexValue_negGates}
\leanok
\uses{BoolCircuit.DAGGate.eval, BoolCircuit.StraightLine.negGates, BoolCircuit.runWith_getD_gate, BoolCircuit.runWith_getD_of_lt, BoolCircuit.vertexValue}
\difficulty{3}
Let $n$ be a natural number, let $x$ be a Boolean assignment to its $n$ inputs, and let
$k$ be an input index. In the gate list consisting of the negation gates for all inputs,
the value at vertex $n+k$ is the negation of $x_k$.
\end{theorem}

\begin{theorem}[Values of straight-line operands]
\lean{BoolCircuit.StraightLine.vertexValue_operand}
\label{BoolCircuit.StraightLine.vertexValue_operand}
\leanok
\uses{BoolCircuit.SLOperand, BoolCircuit.SLOperand.Before, BoolCircuit.SLOperand.eval, BoolCircuit.SLOperand.vertex, BoolCircuit.SLStmt, BoolCircuit.StraightLine.runLines, BoolCircuit.StraightLine.slpGates, BoolCircuit.StraightLine.vertexValue_negGates, BoolCircuit.vertexValue, BoolCircuit.vertexValue_append, BoolCircuit.vertexValue_input}
\difficulty{4}
Let $n$ be a natural number, let $ss$ be a list of straight-line statements on $n$
inputs, and let $x$ be a Boolean assignment to those inputs. Suppose that for every $j$
less than the length of $ss$, the value at vertex $2n+j$ in the gate list for $ss$
equals the value at position $j$ in the execution of $ss$ on $x$ (with an out-of-range
list position interpreted as false). Then every operand $o$ that is available before
$ss$ has the same value in that gate list at its associated vertex as it has when
evaluated on $x$ and the execution values of $ss$.
\end{theorem}

\begin{theorem}[Gate values for valid straight-line lines]
\lean{BoolCircuit.StraightLine.vertexValue_slpGates_line}
\label{BoolCircuit.StraightLine.vertexValue_slpGates_line}
\leanok
\uses{BoolCircuit.SLStmt, BoolCircuit.SLStmt.Before, BoolCircuit.SLStmt.eval, BoolCircuit.SLStmt.toGate_eval, BoolCircuit.StraightLine.LinesValid, BoolCircuit.StraightLine.length_runLines, BoolCircuit.StraightLine.runLines, BoolCircuit.StraightLine.runLines_snoc, BoolCircuit.StraightLine.slpGates, BoolCircuit.StraightLine.slpGates_snoc, BoolCircuit.StraightLine.vertexValue_operand, BoolCircuit.vertexValue, BoolCircuit.vertexValue_append, BoolCircuit.vertexValue_last}
\difficulty{6}
Let $n$ be a natural number, let $ss$ be a list of straight-line statements on $n$
inputs in which each statement reads only input variables, their negations, or earlier
lines, and let $x$ be a Boolean assignment to the inputs. For every $j<|ss|$, the value
at vertex $2n+j$ in the gate list for $ss$ equals the value at position $j$ in the
execution of $ss$ on $x$, with an out-of-range list position interpreted as false.
\end{theorem}

\begin{definition}[Straight-line program circuit]
\lean{BoolCircuit.StraightLineProgram.toCircuit}
\label{BoolCircuit.StraightLineProgram.toCircuit}
\leanok
\uses{BoolCircuit.DAGCircuit, BoolCircuit.GatesAcyclic.snoc, BoolCircuit.StraightLine.length_slpGates, BoolCircuit.StraightLine.slpGates, BoolCircuit.StraightLine.slpGates_acyclic, BoolCircuit.StraightLineProgram.length, BoolCircuit.constGate}
For a straight-line program on $n$ inputs, form a directed acyclic circuit whose gates
are the program's gates in order, appending a constant-false gate when the program has
no statements. Its output vertex is $2n+(m-1)$, where $m$ is the program length.
\end{definition}

\begin{theorem}[Evaluation preserved by circuit translation]
\lean{BoolCircuit.StraightLineProgram.toCircuit_eval}
\label{BoolCircuit.StraightLineProgram.toCircuit_eval}
\leanok
\uses{BoolCircuit.DAGCircuit.eval, BoolCircuit.StraightLine.slpGates, BoolCircuit.StraightLine.vertexValue_slpGates_line, BoolCircuit.StraightLineProgram.eval, BoolCircuit.StraightLineProgram.length, BoolCircuit.StraightLineProgram.toCircuit, BoolCircuit.StraightLineProgram.values, BoolCircuit.vertexValue, BoolCircuit.vertexValue_last}
\difficulty{4}
Let $n$ be a natural number, let $P$ be a straight-line program on $n$ Boolean inputs,
and let $x$ be an assignment of Boolean values to those inputs. Evaluating the circuit
associated with $P$ on $x$ gives the same Boolean value as evaluating $P$ on $x$.
\end{theorem}

\begin{theorem}[Circuit translation has fan-in at most two]
\lean{BoolCircuit.StraightLineProgram.toCircuit_isFaninTwo}
\label{BoolCircuit.StraightLineProgram.toCircuit_isFaninTwo}
\leanok
\uses{BoolCircuit.DAGCircuit.IsFaninTwo, BoolCircuit.SLOp.kind, BoolCircuit.SLStmt.toGate, BoolCircuit.StraightLine.negGates, BoolCircuit.StraightLine.slpGates, BoolCircuit.StraightLineProgram.toCircuit, BoolCircuit.constGate}
\difficulty{5}
Every straight-line program on $n$ inputs translates to a circuit in which each gate has
at most two distinct input vertices, and every negation gate has exactly one input
vertex.
\end{theorem}

\begin{theorem}[Size of the translated circuit]
\lean{BoolCircuit.StraightLineProgram.toCircuit_size}
\label{BoolCircuit.StraightLineProgram.toCircuit_size}
\leanok
\uses{BoolCircuit.DAGCircuit.size, BoolCircuit.StraightLine.length_slpGates, BoolCircuit.StraightLineProgram.length, BoolCircuit.StraightLineProgram.toCircuit}
\difficulty{3}
For every natural number $n$ and every straight-line program $P$ on $n$ inputs, the size
of its translated circuit is $2n+\max(|P|,1)$, where $|P|$ is the number of statements
in $P$.
\end{theorem}

\begin{theorem}[One-line versus size-two circuits for conjunction]
\lean{BoolCircuit.exists_straightLine_and_not_dagCircuit}
\label{BoolCircuit.exists_straightLine_and_not_dagCircuit}
\leanok
\uses{BoolCircuit.DAGCircuit, BoolCircuit.DAGCircuit.eval, BoolCircuit.DAGCircuit.size, BoolCircuit.DAGCircuit.values_getD_input, BoolCircuit.StraightLineProgram, BoolCircuit.StraightLineProgram.eval, BoolCircuit.StraightLineProgram.length}
\difficulty{5}
There is a valid straight-line program of length one that computes conjunction on two
Boolean inputs, but there is no directed acyclic circuit of size at most two that
computes conjunction on those inputs.
\end{theorem}

\begin{definition}[Circuit for exclusive or]
\lean{BoolCircuit.xorCircuit}
\label{BoolCircuit.xorCircuit}
\leanok
\uses{BoolCircuit.DAGCircuit}
The Boolean circuit with inputs $x_0,x_1\in\{0,1\}$ has, in order, gates computing $\neg
x_0$, $\neg x_1$, $(\neg x_0)\land x_1$, $x_0\land(\neg x_1)$, and $((\neg x_0)\land
x_1)\lor(x_0\land(\neg x_1))$. Its output is the value of the fifth gate, so the circuit
computes exclusive or.
\end{definition}

\begin{theorem}[Evaluation of the exclusive-or circuit]
\lean{BoolCircuit.xorCircuit_eval}
\label{BoolCircuit.xorCircuit_eval}
\leanok
\uses{BoolCircuit.DAGCircuit.eval, BoolCircuit.DAGCircuit.values, BoolCircuit.DAGGate.eval, BoolCircuit.runWith, BoolCircuit.xorCircuit}
\difficulty{3}
Let $x$ assign Boolean values to two inputs. The specified Boolean circuit for
exclusive-or evaluates to true exactly when the two input values differ.
\end{theorem}

\begin{theorem}[Exclusive-or circuit has fan-in at most two]
\lean{BoolCircuit.xorCircuit_isFaninTwo}
\label{BoolCircuit.xorCircuit_isFaninTwo}
\leanok
\uses{BoolCircuit.DAGCircuit.IsFaninTwo, BoolCircuit.xorCircuit}
\difficulty{4}
The specified directed acyclic circuit for exclusive-or is well formed, meaning that
each gate's argument list has no repeated indices and every negation gate has exactly
one argument. Each gate has at most two arguments.
\end{theorem}

\begin{theorem}[Size of the XOR circuit]
\lean{BoolCircuit.xorCircuit_size}
\label{BoolCircuit.xorCircuit_size}
\leanok
\uses{BoolCircuit.DAGCircuit.size, BoolCircuit.xorCircuit}
\difficulty{1}
The size of the fixed directed acyclic Boolean circuit for XOR is $7$.
\end{theorem}

\begin{definition}[Straight-line program for exclusive or]
\lean{BoolCircuit.xorProgram}
\label{BoolCircuit.xorProgram}
\leanok
\uses{BoolCircuit.SLOperand.Before, BoolCircuit.SLStmt.Before, BoolCircuit.StraightLineProgram}
For Boolean inputs $x_0$ and $x_1$, the program has three statements in order: the
conjunction $(\neg x_0)\land x_1$, the conjunction $x_0\land(\neg x_1)$, and the
disjunction of the values of the first two statements.
\end{definition}

\begin{theorem}[Evaluation of the XOR program]
\lean{BoolCircuit.xorProgram_eval}
\label{BoolCircuit.xorProgram_eval}
\leanok
\uses{BoolCircuit.SLOp.apply, BoolCircuit.SLOperand.eval, BoolCircuit.SLStmt.eval, BoolCircuit.StraightLine.runLines, BoolCircuit.StraightLineProgram.eval, BoolCircuit.StraightLineProgram.length, BoolCircuit.StraightLineProgram.values, BoolCircuit.xorProgram}
\difficulty{4}
For every assignment $x$ of Boolean values to two variables, the XOR program evaluates
to true if and only if the two assigned values differ.
\end{theorem}

\begin{theorem}[Length of the XOR program]
\lean{BoolCircuit.xorProgram_length}
\label{BoolCircuit.xorProgram_length}
\leanok
\uses{BoolCircuit.StraightLineProgram.length, BoolCircuit.xorProgram}
\difficulty{1}
The XOR straight-line program has length $3$.
\end{theorem}
